\begin{table}[htbp]
\centering
\small
\caption{Bloco 2 (openai/gpt-4o-mini): LLM com 40 exemplos versus classificadores clássicos.}
\label{tab:bloco2_E_LLM_vs_baselines_nmax_gpt4mini}
\begin{tabular}{lrlrrr}
\toprule
 & \multicolumn{5}{c}{Acur.\ vs espec.\ (\%)} \\ \cmidrule(lr){2-6}
Especialista & LLM (mista) & LLM (melhor estratégia) & Logistic Regression & SVM (RBF) & k-NN (k=5) \\
\midrule
E1 & 97,3 & 97,6 (Difícil) & 100,0 & 95,2 & 91,2 \\
E2 & 94,5 & 96,0 (Aleatória) & 100,0 & 95,8 & 93,3 \\
E3 & 93,6 & 97,6 (Difícil) & 100,0 & 99,4 & 98,2 \\
\bottomrule
\multicolumn{6}{p{0.95\linewidth}}{\footnotesize Acurácia (LLM vs especialista), em \%, com n = 40 exemplos (maior n da grade). Os classificadores clássicos (k-NN, Regressão Logística, SVM-RBF) são treinados nos MESMOS n exemplos da estratégia mista e avaliados no MESMO conjunto de teste; não dependem do LLM. Coluna LLM (mista): mesma estratégia dos baselines.} \\
\end{tabular}
\end{table}
